\documentclass[11pt,a4paper]{article}
\usepackage[margin=2.6cm]{geometry}
\usepackage{amsmath,amssymb,amsthm}
\usepackage{tikz}
\usepackage{booktabs}
\usepackage[hidelinks]{hyperref}

\newtheorem*{paperthm}{\paperthmname}
\newcommand{\paperthmname}{}
\theoremstyle{remark}
\newtheorem*{example}{Worked example}
\newtheorem*{checkit}{Check it yourself}
\newcommand{\avg}[1]{\left\langle #1 \right\rangle}
\newcommand{\dd}{\mathrm{d}}
\newcommand{\nb}[1]{\texttt{validation/#1}}

\title{The pizzetti proofs, explained\\[4pt]\large Part 2: The mean-value expansion and Appendix B}
\author{Companion to the pizzetti paper (Section 4.2 and Appendix B, including B.1)}
\date{}

\begin{document}
\maketitle

\begin{abstract}
\noindent Part 1 grouped the series by powers of $\varepsilon$. Here the same series is regrouped
by powers of the planet's size $p$. This connects it to a 1909 formula by Pizzetti for the
average of a function over a disk, and to the small-planet and qpower2 approximations used in
practice. We then prove the two error bounds of Appendix~B, going through the proof of
Appendix~B.1 line by line. The tools are Taylor expansion, the Laplacian, Gauss's divergence
theorem in two dimensions (familiar from electrostatics), and one theorem about power series
(Abel's). Every formula is checked in \nb{03\_meanvalue\_appendixB.ipynb}.
\end{abstract}

\tableofcontents

%% ===================================================================
\section{The idea: average of a function over a small disk}

Take any smooth function $f(x,y)$ and average it over a disk of radius $p$ centred at a point $c$.
If $p$ is small, the average is close to $f(c)$. How close? Expand $f$ around $c$ in a Taylor
series. Terms of odd order (like $x$, $x^3$, $xy^2$) average to zero by symmetry: for every point
of the disk the mirror point gives the opposite value. Terms of second order give
$\avg{x^2}=\avg{y^2}=p^2/4$ and $\avg{xy}=0$, so the second-order part of the Taylor series averages
to
\[
  \tfrac12 f_{xx}\avg{x^2}+\tfrac12 f_{yy}\avg{y^2}=\frac{p^2}{8}\,(f_{xx}+f_{yy})=\frac{p^2}{8}\,\Delta f(c),
\]
with the Laplacian $\Delta f=f_{xx}+f_{yy}$. Continuing to all orders gives the
\emph{mean-value expansion}
\begin{equation}
  \avg{f}=\sum_{n=0}^\infty\frac{(p^2/4)^n}{n!\,(n+1)!}\,\Delta^nf(c)
  =f(c)+\frac{p^2}{8}\Delta f(c)+\frac{p^4}{192}\Delta^2f(c)+\dots
  \label{eq:mv}
\end{equation}
($\Delta^n$ is the Laplacian applied $n$ times). Only Laplacians appear because averaging over a
disk does not prefer any direction, and the Laplacian is the only second-order operation without
a preferred direction. The formula is the area version of a formula Pizzetti (1909) gave for
averages over spheres; Section~\ref{sec:proof} proves it, with an exact remainder.

\begin{example}
For $f=x^2+y^2$ and $c=0$: the disk average of $r^2$ is $p^2/2$, and \eqref{eq:mv} gives
$0+\frac{p^2}{8}\cdot4=p^2/2$, since $\Delta(x^2+y^2)=4$ and all higher Laplacians vanish.
\end{example}

%% ===================================================================
\section{Applied to a star: the familiar approximations}

For the star, $f=(1-r^2)^\gamma$, the intensity term $\mu^{2\gamma}$, and $c$ is the planet's
centre at distance $z$. The first term of \eqref{eq:mv}, $f(c)=(1-z^2)^\gamma$, is the
\emph{small-planet approximation}: the brightness at the planet's centre times its area. The first
two terms, $I+\frac{p^2}{8}\Delta I$, are exactly the qpower2 approximation of Maxted \& Gill
(2019) for the power-2 law, and the first-order Rossiter--McLaughlin formula of Ohta et al.\ (2005).
So these known approximations are the beginning of one series, and the paper's polynomial is the
same series, summed further.

\paragraph{Computing the Laplacians.} For a function of $r$ alone, the Laplacian in the plane is
$f''+f'/r$. For a power of $w=1-r^2$ this gives (paper, Eq.~22)
\begin{equation}
  \Delta\,w^{s}=4s(s-1)\,w^{s-2}-4s^2\,w^{s-1}.
  \label{eq:lap}
\end{equation}
\emph{Derivation:} $\frac{\dd}{\dd r}w^s=-2rs\,w^{s-1}$ and
$\frac{\dd^2}{\dd r^2}w^s=4r^2s(s-1)w^{s-2}-2s\,w^{s-1}$; add $\frac1r\frac{\dd}{\dd r}w^s=-2s\,w^{s-1}$
and use $r^2=1-w$. Applying \eqref{eq:lap} repeatedly produces every $\Delta^nf$ as a short sum of
powers of $w$. For $\gamma=\tfrac12$ the first four terms of \eqref{eq:mv} become the paper's
Eq.~(23):
\[
  D_{1/2}=\sqrt{w}\Big[1-\frac{p^2}{8}(v+v^2)+\frac{p^4}{192}(v^2+6v^3-15v^4)
  -\frac{p^6}{1024}(v^3+15v^4-105v^5+105v^6)+\dots\Big],\quad v=\frac1w .
\]

\paragraph{Why this is the same series as in Part 1.} Part 1 wrote $D_\gamma$ as a sum over
powers of $\varepsilon$. Each term there is a polynomial in $p^2$ and $z^2$. Collecting all terms
with the same power $p^{2n}$ gives a term $T_n$ (paper, Eq.~20). Both series are power series in
$p^2$ for the same quantity, and a power series has only one set of coefficients. So $T_n$
\emph{must} equal the $n$-th term of \eqref{eq:mv}. (Regrouping is allowed because the series of
Part 1 converges absolutely.)

\begin{example}
For $\gamma=\tfrac12$, $p=0.1$, $z=0.5$: $T_0=0.866025$, $T_1=-3.37\times10^{-3}$,
$T_2=-1.42\times10^{-5}$, $T_3=-1.67\times10^{-7}$. All terms after the first are negative, like
the terms of Part 1. (They are sums of Part 1's terms, which all have one sign for $0<\gamma<1$.)
\end{example}

\begin{checkit}
\nb{03\_meanvalue\_appendixB.ipynb}, cells 3a--3b: \eqref{eq:lap} and the $\gamma=\tfrac12$
expansion are verified symbolically (sympy computes the Laplacians from scratch), and $T_n$ agrees
with the $n$-th mean-value term to 40 digits for several $\gamma$.
\end{checkit}

%% ===================================================================
\section{Appendix B.1, step by step: the mean-value expansion with remainder}
\label{sec:proof}

We now prove \eqref{eq:mv} again, this time keeping track of what is left over when we stop after
$N$ terms. The argument has four steps.

\subsection*{Step 1: Circle averages and Gauss's law}

Let $m_f(\rho)$ be the average of $f$ over the \emph{circle} of radius $\rho$ around $c$. How does
it change with $\rho$? Its derivative is the average of the outward slope $\partial f/\partial\rho$
over the circle. The integral of the outward slope around the circle is the \emph{flux} of the
gradient $\nabla f$ through the circle, and by the divergence theorem (Gauss's law in two
dimensions) that flux equals the integral of $\nabla\cdot\nabla f=\Delta f$ over the enclosed disk:
\[
  2\pi\rho\,m_f'(\rho)=\oint\frac{\partial f}{\partial\rho}\,\dd\ell
  =\int_{\text{disk of radius }\rho}\Delta f\,\dd A
  =2\pi\int_0^\rho s\,m_{\Delta f}(s)\,\dd s .
\]
In the last step the disk is cut into thin rings of radius $s$ and circumference $2\pi s$, on each
of which the average of $\Delta f$ is $m_{\Delta f}(s)$. In words: \emph{the circle average grows at
a rate set by the average Laplacian inside it}, just as the electric field through a circle is set
by the enclosed charge.

\subsection*{Step 2: Integrate once}

Divide by $2\pi\rho$, integrate from $0$ to $\rho$, and use $m_f(0)=f(c)$:
\[
  m_f(\rho)=f(c)+\int_0^\rho\frac{1}{r}\int_0^r s\,m_{\Delta f}(s)\,\dd s\,\dd r .
\]
Swap the two integrals (the region is $0\le s\le r\le\rho$; for fixed $s$, $r$ runs from $s$ to
$\rho$, and $\int_s^\rho\dd r/r=\ln(\rho/s)$):
\begin{equation}
  m_f(\rho)=f(c)+\int_0^\rho s\ln\frac{\rho}{s}\;m_{\Delta f}(s)\,\dd s
  =f(c)+(\mathcal T m_{\Delta f})(\rho).
  \label{eq:T}
\end{equation}
We call the integral operation $\mathcal T$. Two properties matter:
\begin{itemize}
  \item its weight $s\ln(\rho/s)$ is \emph{never negative} for $0\le s\le\rho$;
  \item it maps powers to powers, $\mathcal T[\rho^{2j}]=\rho^{2j+2}/(2j+2)^2$;
        to see this, substitute $s=\rho u$ and use
        $\int_0^1u^{2j+1}\ln(1/u)\,\dd u=1/(2j+2)^2$.
\end{itemize}

\subsection*{Step 3: Repeat}

Equation \eqref{eq:T} holds for any smooth function, so apply it to $\Delta f$ inside the integral:
$m_{\Delta f}=\Delta f(c)+\mathcal T m_{\Delta^2f}$. Substituting, and using
$\mathcal T[1]=\rho^2/4$,
\[
  m_f(\rho)=f(c)+\frac{\rho^2}{4}\Delta f(c)+\mathcal T^2m_{\Delta^2f}(\rho).
\]
Repeating $N+1$ times gives
\begin{equation}
  m_f(\rho)=\sum_{j=0}^{N}\frac{\rho^{2j}}{4^j(j!)^2}\,\Delta^jf(c)
  +\underbrace{\mathcal T^{N+1}m_{\Delta^{N+1}f}(\rho)}_{\text{remainder}} .
  \label{eq:circle}
\end{equation}
The coefficients follow from the power rule: $\mathcal T$ applied $j$ times to $1$ gives
$\rho^{2j}/(2^2\cdot4^2\cdots(2j)^2)=\rho^{2j}/(4^j(j!)^2)$.

\subsection*{Step 4: From circles to the disk, and the size of the remainder}

The disk average is the circle averages weighted by circumference:
$\avg{f}=\frac{2}{p^2}\int_0^p\rho\,m_f(\rho)\,\dd\rho$. Applied to the terms of \eqref{eq:circle},
$\frac{2}{p^2}\int_0^p\rho\cdot\rho^{2j}\,\dd\rho=\frac{p^{2j}}{j+1}$, which turns
$\frac{1}{4^j(j!)^2}$ into $\frac{1}{4^j\,j!\,(j+1)!}$: exactly the coefficients of \eqref{eq:mv}.

The remainder is built from $m_{\Delta^{N+1}f}$ by repeated application of $\mathcal T$ (weights
$\ge0$) and the final disk average (weights $\ge0$). So it is a \emph{weighted average with
non-negative weights} of the circle averages of $\Delta^{N+1}f$:
\[
  \text{remainder}=\int_0^pK(s)\,m_{\Delta^{N+1}f}(s)\,\dd s,\qquad K(s)\ge0 .
\]
Two consequences:
\begin{enumerate}
  \item The total weight $\int K$ is what the remainder would be if $\Delta^{N+1}f$ were the
        constant 1. Then the next term of the series is the whole remainder, so
        $\int K=(p^2/4)^{N+1}/((N+1)!\,(N+2)!)$.
  \item Each circle average lies between the smallest and largest value of $\Delta^{N+1}f$ on the
        disk, and so does any weighted average of them. Hence there is a point $\xi$ of the disk
        with
\end{enumerate}
\begin{equation}
  \avg{f}-\sum_{n=0}^{N}\frac{(p^2/4)^n}{n!(n+1)!}\Delta^nf(c)
  =\frac{(p^2/4)^{N+1}}{(N+1)!\,(N+2)!}\,\Delta^{N+1}f(\xi).
  \label{eq:lagrange}
\end{equation}
This is the analogue of the Lagrange remainder of a Taylor series: the next term, evaluated at an
unknown point instead of the centre.

%% ===================================================================
\section{Proposition 1: a two-sided bound}

\renewcommand{\paperthmname}{Proposition 1 (as in the paper)}
\begin{paperthm}
For $f=(1-r^2)^\gamma$ with $0<\gamma<1$, all $\Delta^nf$ with $n\ge1$ are negative and their size
increases with $r$. The remainder of the mean-value expansion after $N$ terms therefore lies between
the right-hand side of \eqref{eq:lagrange} evaluated at $r=\max(0,z-p)$ and at $r=z+p$.
\end{paperthm}

\paragraph{Why all Laplacians are negative.} Write $f$ as a power series in $r^2$:
$f=\sum_m\beta_mr^{2m}$ with the binomial coefficients $\beta_m$ of Part 1, all negative for
$m\ge1$ when $0<\gamma<1$. In the plane, $\Delta r^{k}=k^2r^{k-2}$, so
$\Delta^nr^{2m}=4^n\bigl(m!/(m-n)!\bigr)^2r^{2(m-n)}$ (positive). Hence
\[
  \Delta^nf=\sum_{m\ge n}\beta_m\,4^n\Bigl(\frac{m!}{(m-n)!}\Bigr)^2r^{2(m-n)}
\]
is a power series in $r^2$ with only negative coefficients (for $n\ge1$). Such a function is
negative, and its size grows with $r$.

\paragraph{Where on the disk.} The unknown point $\xi$ lies somewhere on the planet's disk. The
smallest distance from the stellar centre on that disk is $z-p$ if the planet does not cover the
centre ($z\ge p$), and $0$ if it does ($z<p$); the largest is $z+p$. Since $|\Delta^{N+1}f|$
increases with $r$, the remainder lies between its values at these two radii.

\paragraph{A correction found by the checks.} An earlier draft of the paper wrote $|z-p|$ for the
smallest radius. For $z<p$ that is wrong: the disk then contains the stellar centre, where $r=0$.
At $z=0$ the two ends of the old bound coincide, and the true remainder is not equal to that
single value. The validation notebook found this, and the paper now has $\max(0,z-p)$.

\begin{example}
$\gamma=\tfrac12$, $p=0.1$, $z=0.5$:
\[
\begin{array}{c|ccc}
 N & \text{lower end} & \text{true remainder} & \text{upper end}\\\hline
 0 & -4.00\times10^{-3} & -3.38\times10^{-3} & -2.99\times10^{-3}\\
 1 & -2.67\times10^{-5} & -1.43\times10^{-5} & -8.87\times10^{-6}\\
 2 & -5.03\times10^{-7} & -1.70\times10^{-7} & -7.13\times10^{-8}
\end{array}
\]
And at $z=0$, $N=1$: the true remainder is $-4.182\times10^{-6}$; the corrected bound gives
$[-4.359,\,-4.167]\times10^{-6}$ (contains it); the old one gave the single value
$-4.359\times10^{-6}$ (does not).
\end{example}

\paragraph{Limitation.} The lower end uses the largest radius $z+p$, where $\Delta^{N+1}f$ is
largest. Near the stellar edge it blows up, roughly like $(p/(1-z-p))^{2N+2}$: for $N=2$ the bound
spans $[-5\times10^{-7},-7\times10^{-8}]$ at $z=0.5$ but $[-5.1,\,-1\times10^{-5}]$ at $z=0.88$
($p=0.1$). Proposition~1 is therefore useful only away from the edge. Proposition~2 fixes this for
the quadratic law.

%% ===================================================================
\section{Proposition 2: a bound valid up to the edge (Abel)}

\renewcommand{\paperthmname}{Proposition 2 (as in the paper)}
\begin{paperthm}
For $\gamma=\tfrac12$, $t=p^2$, $t_1=(1-z)^2$ and $0\le p\le1-z$,
\[
  0\le S^{\rm P}_N(z,p)-D_{1/2}(z,p)\le\Bigl(\frac{p}{1-z}\Bigr)^{2N+2}
  \Bigl[S^{\rm P}_N(z,1-z)-D_{1/2}(z,1-z)\Bigr],
\]
where $S^{\rm P}_N=\sum_{n\le N}T_n$, and $D_{1/2}(z,1-z)$, the average over the planet just
touching the edge from inside, has the elementary formula of the paper's Eq.~(B4).
\end{paperthm}

The idea: view $D_{1/2}$ at fixed $z$ as a power series in $t=p^2$,
$D_{1/2}=\sum_nd_nt^n$ (so $T_n=d_nt^n$). The largest planet that fits is the one touching the
edge, $t=t_1$. If we knew the error at $t_1$, we could scale it down to smaller planets. Four small
steps:

\begin{enumerate}
  \item \emph{All $d_n$ with $n\ge1$ are negative} (the $T_n$ are, see above). So the remainder
        $D-S^{\rm P}_N=\sum_{n>N}d_nt^n$ is negative: $S^{\rm P}_N\ge D$ for every $t<t_1$.
        That is the left inequality.
  \item \emph{The series still converges at the edge, $t=t_1$.} The partial sums at $t_1$ decrease
        as $N$ grows (we keep adding negative terms). They cannot fall below $D(z,1-z)$: for
        $t<t_1$ we have $S^{\rm P}_N\ge D$, and letting $t\to t_1$ keeps the inequality because
        both sides are continuous. A decreasing sequence that is bounded below converges.
  \item \emph{And it converges to the right value.} Abel's theorem says: if a power series converges
        at the end point of its interval, its sum there is the limit of its values inside. So
        $\sum_nd_nt_1^n=D(z,1-z)$, and the error at the edge is
        $\sum_{n>N}|d_n|t_1^n=S^{\rm P}_N(z,1-z)-D(z,1-z)$.
  \item \emph{Scale down.} For $t<t_1$ and $n>N$, $t^n=t_1^n\,(t/t_1)^n\le t_1^n\,(t/t_1)^{N+1}$.
        Summing, $\sum_{n>N}|d_n|t^n\le(t/t_1)^{N+1}\sum_{n>N}|d_n|t_1^n$, and
        $(t/t_1)^{N+1}=(p/(1-z))^{2N+2}$. That is the right inequality.
\end{enumerate}

\begin{example}
$\gamma=\tfrac12$, $p=0.1$, $z=0.5$: the true errors $S^{\rm P}_N-D$ are $3.38\times10^{-3}$,
$1.43\times10^{-5}$ and $1.70\times10^{-7}$ for $N=0,1,2$; the Abel bounds are
$3.94\times10^{-3}$, $2.30\times10^{-5}$ and $3.54\times10^{-7}$. The tangent-disk formula gives
$D_{1/2}(0.5,0.5)=0.767449091228817$, the same as the defining integral to all digits shown.
\end{example}

\begin{checkit}
\nb{03\_meanvalue\_appendixB.ipynb}, cells 3c--3e: the tangent-disk formula against the defining
integral (agreement $10^{-36}$); Proposition~1 with both the corrected and the old smallest radius
(only the old one fails, and only for $z<p$); Proposition~2 with no violations.
\end{checkit}

%% ===================================================================
\section{Summary}

Regrouping the series of Part 1 by powers of the planet size gives Pizzetti's mean-value expansion,
whose first one and two terms are the small-planet and qpower2 approximations. Appendix~B.1 proves
the expansion with an exact remainder using one physical idea, Gauss's law for the gradient of $f$,
and the fact that averages with non-negative weights stay between the minimum and maximum.
Proposition~1 turns this into a two-sided bound that works away from the stellar edge;
Proposition~2 uses the sign of the coefficients and Abel's theorem to obtain a bound that works up
to the edge. Neither is used by the pizzetti code, which relies on Theorem~2 (Part 1); they explain
how the known approximations fit in and how accurate they are.

\medskip\noindent\textbf{Next:} Part 3, the phases in which the planet crosses the stellar edge
(Sections 5.1 and 5.3 and Appendix C).

\end{document}
